\documentclass[11pt]{article}

\usepackage[utf8]{inputenc}
\usepackage[T1]{fontenc}
\usepackage{lmodern}
\usepackage{geometry}
\usepackage{amsmath,amssymb,amsfonts,amsthm,mathtools}
\usepackage{bm}
\usepackage{enumitem}
\usepackage{microtype}
\usepackage{hyperref}
\usepackage{titlesec}
\usepackage{booktabs}
\usepackage{array}
\usepackage{setspace}
\usepackage{mathrsfs}
\usepackage{physics}

\geometry{margin=1in}
\hypersetup{colorlinks=true, linkcolor=black, urlcolor=blue, citecolor=black}
\setlist[enumerate]{topsep=0.4em,itemsep=0.35em}
\setlist[itemize]{topsep=0.3em,itemsep=0.25em}
\titleformat{\section}{\large\bfseries}{\thesection.}{0.5em}{}
\titleformat{\subsection}{\normalsize\bfseries}{\thesubsection.}{0.5em}{}
\setstretch{1.06}

\newcommand{\Aether}{\AE{}ther}
\newcommand{\Aflow}{\AE{}ther-flow}
\newcommand{\TheoryName}{\Aflow{} Interpretation of Relativity}
\newcommand{\TheoryProgram}{\Aether{} / \Aflow{} framework}
\newcommand{\Sub}{\mathcal A}
\newcommand{\ObserverMap}{\Xi}
\newcommand{\BridgeMetric}{\mathfrak g}
\newcommand{\ObserverMetric}{g^{\mathrm{br}}}
\newcommand{\CandidateMetric}{g^{\mathrm{cand}}}
\newcommand{\CompletedMetric}{g^{\sharp}}
\newcommand{\BaseShift}{\Sigma^{\mathrm{IER}}}
\newcommand{\ResponseShift}{\Sigma^{\mathrm{resp}}}
\newcommand{\CompShift}{\Sigma^{\mathrm{cmp}}}
\newcommand{\BaseLift}{\widehat\Sigma^{\mathrm{IER}}}
\newcommand{\ResponseLift}{\widehat\Sigma^{\mathrm{resp}}}
\newcommand{\CompLift}{\widehat\Sigma^{\mathrm{cmp}}}
\newcommand{\ReLongFour}{\widetilde{\mathcal R}_{\mu\nu}^{(\chi,\ge 4)}}
\newcommand{\ResidualCore}{\mathcal V_{\mu\nu}^{\mathrm{res}}}
\newcommand{\ResidualDec}{\mathcal D_{\mu\nu}^{\mathrm{dec}}}
\newcommand{\ResidualLift}{\widehat{\mathcal V}^{\mathrm{res}}}
\newcommand{\CompFunctional}{\mathcal W_{\mathrm{cmp}}}
\newcommand{\CompAction}{S_{\mathrm{cmp}}}
\newcommand{\CompletedAction}{S_{\mathrm{cand}}^{\sharp}}
\newcommand{\ChargeFields}{Q_{\mathrm{cp}}}
\newcommand{\SubEinLin}{\widehat{\mathcal E}}
\newcommand{\SubGaugeVec}{\widehat{\mathcal B}}
\newcommand{\SubGaugeBook}{\widehat{\mathcal G}}
\newcommand{\SubReducedEin}{\widehat{\mathcal L}}
\newcommand{\FreeProj}{\Pi_{\mathrm{free}}}
\newcommand{\GaugeAux}{Y}
\newcommand{\ModeAux}{Z}
\newcommand{\ReducedEin}{\mathcal L_{\ObserverMetric}}
\newcommand{\GaugeBook}{\mathcal G_{\ObserverMetric}}
\newcommand{\EinLin}{\mathcal E_{\ObserverMetric}}
\newcommand{\EinQuad}{\mathcal Q_{\ge 2}^{\mathrm{Ein}}}
\newcommand{\CompMix}{\mathcal Q_{\mu\nu}^{\mathrm{cmp,mix}}}
\newcommand{\CompQuad}{\mathcal Q_{\mu\nu}^{\mathrm{cmp}}}
\newcommand{\CompDec}{\mathcal D_{\mu\nu}^{\sharp}}
\newcommand{\DeltaTComp}{\Delta T_{\mu\nu}^{\mathrm{cmp}}}
\newcommand{\MatterAct}{S_{\mathrm{matter}}}
\newcommand{\epsIR}{\varepsilon}

\newtheorem{definition}{Definition}
\newtheorem{proposition}{Proposition}
\newtheorem{remark}{Remark}
\newtheorem{corollary}{Corollary}
\newtheorem{theorem}{Theorem}

\title{\textbf{The \TheoryName{}}\\[0.4em]
\large Substrate-Response / Self-Response Charge-Polarization Candidate\\
\large Quartic-Completion Substrate-Action Note}
\author{}
\date{}

\begin{document}

\maketitle

\begin{abstract}
The quartic-completion substrate-anchoring note already identifies the lifted quartic self-response source \(\ResidualLift\), introduces a reconstruction-level completion functional, and shows that its elimination reproduces the observer-side solve
\[
\ReducedEin(\CompShift)_{\mu\nu}=-\ResidualCore{}_{\mu\nu}.
\]
The remaining burden is narrower and deeper: derive that reconstruction-level functional from an explicit completion-sector action inside the substrate action class already fixed by the substrate-kinematics manuscript.

The present note performs that step on the same one-metric charge-polarization line. Starting from the candidate substrate action
\[
S_{\mathrm{cand}}
=
\frac{c^3}{16\pi G_N}\int_{\Sub}\sqrt{G}\,\bigl(R[\BridgeMetric]-2\Lambda_{\Sub}\bigr)
+\int_{\Sub}\sqrt{G}\,\mathcal L_{\mathrm{aux}}^{(0)}
+\MatterAct[\ObserverMap^*\BridgeMetric,\psi],
\]
it refines the extra substrate block \(Q\) by adjoining an observer-compatible completion tensor \(H_{AB}\), a reduced-gauge auxiliary covector \(\GaugeAux_A\), and a free-homogeneous-mode multiplier \(\ModeAux_{AB}\). The resulting completion action is
\[
\CompAction[H,\GaugeAux,\ModeAux]
=
\int_{\Sub}\sqrt{G}\,
\left[
\frac12 H^{AB}\SubEinLin(H)_{AB}
-\GaugeAux^A\SubGaugeVec(H)_A
+\frac12 \GaugeAux^A\GaugeAux_A
+(\ResidualLift)^{AB}H_{AB}
+\ModeAux^{AB}(\FreeProj H)_{AB}
\right].
\]
Variation of \(\GaugeAux_A\) derives the reduced operator
\[
\SubReducedEin=\SubEinLin-\SubGaugeBook,
\]
variation of \(\ModeAux_{AB}\) derives the zero-extra-homogeneous-mode condition
\[
\FreeProj H=0,
\]
and variation of \(H_{AB}\) yields the on-shell equation
\[
\SubReducedEin(\CompLift)_{AB}=-\ResidualLift_{AB}.
\]
After observer pullback, the same quartic-completion solve \(\ReducedEin(\CompShift)=-\ResidualCore\) is recovered, so the exact observer-equation closure, the one operative metric, and the recorded Phase D / E and Phase F package are unchanged.

The positive result is disciplined. This note gives an explicit action-level origin of the completion sector inside the existing substrate action class. It does not yet claim a unique ultraviolet completion, a broader finite-amplitude theorem, or promotion beyond the current claim boundary.
\end{abstract}

\section{Introduction}

The active charge-polarization sequence has already fixed the concrete candidate branch, the full Phase D / E infrared-spectrum and universal-coupling verdict, the Phase F controlled GR-limit theorem, the residual-sector verdict, the quartic-completion redesign, and the first substrate-anchoring layer of that redesign. The next honest burden is therefore no longer another packaging pass, no longer stronger promotion language, and no longer a default return to the anchored positive-pair side line. The next burden is to push the quartic-completion sector one step deeper inside the same one-metric charge-polarization branch.

That is the purpose of the present note. The anchoring note already did something important: it showed that the observer-level completion solve can be viewed as the elimination of a reconstruction-level substrate functional. What it did not yet do was derive that functional from an explicit completion-sector action inside the substrate action class already fixed by the derivational-bridge manuscript.

\paragraph{Framework and claim boundary.}
\TheoryProgram{} denotes the broader \Aether{} / \Aflow{} framework built on the claims that the \Aether{} is the underlying four-dimensional substrate of reality, the \Aflow{} is its intrinsic ordered motion, observed three-dimensional space is the local experiential slice of that deeper substrate, S-time is the experienced order of change arising from matter, light, and the \Aflow{}, and observed expansion is the three-dimensional appearance of deeper four-dimensional ordered motion. Within that framework, \TheoryName{} denotes the exact-closure theory in which the effective gravitational dynamics are adopted to be exactly Einsteinian with universal matter coupling. In the active sequence, \emph{adoption} means use of that established relativistic dynamics without claiming substrate derivation, while \emph{derivation} is reserved for a first-principles recovery from explicit substrate variables.

The present note stays strictly inside that boundary. It does not claim a unique microscopic origin of the completion kernel. It does not claim that the full substrate action has now been derived from still deeper microphysics. It does make one narrower positive move: it writes an explicit completion-sector action inside the already-fixed substrate action class, derives from that sector the source coupling, the reduced completion operator, and the zero-extra-homogeneous-mode condition, and then re-shows that elimination yields the same \(\CompShift\) solve already used in the redesign and anchoring notes.

\section{Fixed Starting Point from the Active Sequence}

\begin{definition}[Candidate substrate action class]
\label{def:candidate}
The starting point is the substrate action class already fixed in the substrate-kinematics manuscript:
\begin{equation}
S_{\mathrm{cand}}
=
\frac{c^3}{16\pi G_N}
\int_{\Sub} d^4X\,\sqrt{G}\,\bigl(R[\BridgeMetric]-2\Lambda_{\Sub}\bigr)
+
\int_{\Sub} d^4X\,\sqrt{G}\,\mathcal L_{\mathrm{aux}}^{(0)}[G,U,S,Q]
+
\MatterAct[\ObserverMap^*\BridgeMetric,\psi],
\label{eq:Scand}
\end{equation}
where \(\BridgeMetric_{AB}=G_{AB}-2U_AU_B\), the observer map satisfies \(\ObserverMap^*\BridgeMetric=\ObserverMetric\), and \(\mathcal L_{\mathrm{aux}}^{(0)}\) denotes the auxiliary substrate block already fixed there: unit-flow constraint, order-scalar terms, the original \(Q\)-sector kinetic and mass terms, and the derivative terms controlling the ordered-flow field.
\end{definition}

\begin{remark}
The present note does not alter the bridge metric, the universal matter route, or the role of the Einstein--Hilbert bridge term in \eqref{eq:Scand}. It refines only the extra substrate block \(Q\), and it does so in the specific branch already selected by the charge-polarization line.
\end{remark}

\begin{definition}[Fixed charge-polarization branch]
\label{def:fixed}
The already-recorded charge-polarization branch is treated as fixed in the form
\begin{equation}
\CandidateMetric
=
\ObserverMetric+\BaseShift+\ResponseShift,
\label{eq:gcand}
\end{equation}
with substrate lifts \(\BaseLift_{AB}\) and \(\ResponseLift_{AB}\) satisfying
\begin{equation}
\ObserverMap^*\BaseLift_{AB}=\BaseShift_{\mu\nu},
\qquad
\ObserverMap^*\ResponseLift_{AB}=\ResponseShift_{\mu\nu},
\label{eq:lifts}
\end{equation}
and candidate bridge metric
\begin{equation}
\BridgeMetric_{AB}^{\mathrm{cand}}
=
\BridgeMetric_{AB}+\BaseLift_{AB}+\ResponseLift_{AB}.
\label{eq:bridgecand}
\end{equation}
The already-fixed order facts are
\begin{equation}
\BaseShift=\mathcal O_{\ge 2},
\qquad
\ResponseShift=\mathcal O_{\ge 4},
\label{eq:orders}
\end{equation}
and on every controlled infrared family,
\begin{equation}
\BaseShift=\mathcal O(\epsIR^2),
\qquad
\ResponseShift=\mathcal O(\epsIR^4).
\label{eq:irorders}
\end{equation}
\end{definition}

\begin{definition}[Lifted quartic self-response source]
\label{def:source}
The unresolved observer-side quartic tensor is
\begin{equation}
\ResidualCore
\equiv
\ReLongFour
+\EinQuad[\BaseShift]
+\EinLin(\ResponseShift),
\label{eq:vres}
\end{equation}
and its substrate lift is the explicit tensor
\begin{equation}
\ResidualLift_{AB}
\equiv
\widehat{\mathcal R}_{AB}^{(\chi,\ge 4)}
+\widehat{\mathcal Q}_{AB}^{\mathrm{Ein}}[\BaseLift]
+\widehat{\mathcal E}_{AB}(\ResponseLift),
\label{eq:liftedsource}
\end{equation}
chosen so that
\begin{equation}
\ObserverMap^*\ResidualLift_{AB}
=
\ResidualCore{}_{\mu\nu}.
\label{eq:sourcepull}
\end{equation}
The order bookkeeping is
\begin{equation}
\ResidualLift=\mathcal O_{\ge 4},
\qquad
\ResidualCore=\mathcal O_{\ge 4},
\label{eq:sourceorder}
\end{equation}
and on every controlled infrared family,
\begin{equation}
\ResidualLift=\mathcal O(\epsIR^4),
\qquad
\ResidualCore=\mathcal O(\epsIR^4).
\label{eq:sourceIR}
\end{equation}
\end{definition}

\begin{remark}
Definition \ref{def:source} is the exact reconstruction-level input carried over from the anchoring note. The present manuscript does not change that source. It derives an explicit action-level completion sector whose elimination answers to that same source.
\end{remark}

\section{Completion Sector Inside the Substrate Action Class}

\begin{definition}[Completion-sector refinement of the \(Q\)-block]
\label{def:Qrefine}
Inside the candidate action class \eqref{eq:Scand}, refine the extra substrate block by decomposing
\begin{equation}
Q
=
\bigl(\ChargeFields,\,
H_{AB},\,
\GaugeAux_A,\,
\ModeAux_{AB}\bigr),
\label{eq:Qsplit}
\end{equation}
where:
\begin{enumerate}
    \item \(\ChargeFields\) denotes the already-fixed charge-polarization fields whose elimination produced \(\BaseLift\) and \(\ResponseLift\);
    \item \(H_{AB}\) is an observer-compatible symmetric completion tensor on \(\Sub\);
    \item \(\GaugeAux_A\) is an auxiliary covector used to derive the reduced completion operator;
    \item \(\ModeAux_{AB}\) is a multiplier imposing the vanishing of any extra free homogeneous completion mode.
\end{enumerate}
\end{definition}

\begin{definition}[Linear bridge-response operators]
\label{def:operators}
Let \(\widehat\nabla\) denote the Levi-Civita connection of the unshifted bridge metric \(\BridgeMetric_{AB}\). Let \(\SubEinLin\) be the formally self-adjoint linear bridge-response operator obtained from the quadratic expansion of the Einstein--Hilbert bridge term in \eqref{eq:Scand} in the completion direction \(H_{AB}\):
\begin{equation}
\frac{c^3}{16\pi G_N}
\int_{\Sub}\sqrt{G}\,R[\BridgeMetric+sH]
=
\mathrm{const}
+s\,\mathfrak J[H]
+\frac{s^2}{2}\int_{\Sub}\sqrt{G}\,H^{AB}\SubEinLin(H)_{AB}
+\mathcal O(s^3).
\label{eq:quadEH}
\end{equation}
Define the substrate reduced-gauge vector and gauge-bookkeeping tensor by
\begin{equation}
\SubGaugeVec(H)_A
\equiv
\widehat\nabla^B\!\left(
H_{AB}
-\frac12 \BridgeMetric_{AB}\,\BridgeMetric^{CD}H_{CD}
\right),
\label{eq:subB}
\end{equation}
\begin{equation}
\SubGaugeBook(H)_{AB}
\equiv
\widehat\nabla_{(A}\SubGaugeVec(H)_{B)}
-\frac12 \BridgeMetric_{AB}\widehat\nabla^C\SubGaugeVec(H)_C,
\label{eq:subG}
\end{equation}
and the reduced completion operator
\begin{equation}
\SubReducedEin(H)_{AB}
\equiv
\SubEinLin(H)_{AB}-\SubGaugeBook(H)_{AB}.
\label{eq:subL}
\end{equation}
These operators are chosen compatibly with the observer map so that
\begin{equation}
\ObserverMap^*\bigl(\SubEinLin(H)_{AB}\bigr)
=
\EinLin(\ObserverMap^*H)_{\mu\nu},
\qquad
\ObserverMap^*\bigl(\SubGaugeBook(H)_{AB}\bigr)
=
\GaugeBook(\ObserverMap^*H)_{\mu\nu},
\label{eq:intertwiners}
\end{equation}
hence
\begin{equation}
\ObserverMap^*\bigl(\SubReducedEin(H)_{AB}\bigr)
=
\ReducedEin(\ObserverMap^*H)_{\mu\nu}.
\label{eq:intertwineL}
\end{equation}
\end{definition}

\begin{remark}
The reduced operator \(\SubReducedEin\) is therefore not postulated independently. Its unreduced part comes from the same Einstein--Hilbert bridge term already fixed in the substrate action class, while its gauge-bookkeeping part is produced by the auxiliary covector \(\GaugeAux_A\) introduced below.
\end{remark}

\begin{definition}[Free homogeneous completion projector]
\label{def:freeproj}
Let \(\FreeProj\) denote the projector extracting the \emph{free homogeneous completion data} in the same reduced-harmonic Coulomb-plus-regular completion class used in the redesign and anchoring notes. By construction:
\begin{enumerate}
    \item \(\FreeProj H\) records only the independent homogeneous completion mode after the forced Coulomb-plus-regular particular branch has been separated;
    \item \(\FreeProj\circ\SubReducedEin=0\);
    \item \(\FreeProj\ResidualLift=0\).
\end{enumerate}
The third statement is the action-level form of the earlier admissibility rule: the quartic self-response source determines only the forced completion solve and carries no independent free homogeneous completion data.
\end{definition}

\begin{proposition}[Derived source coupling from the unresolved branch expansion]
\label{prop:sourcecoupling}
When the candidate action class is expanded on the fixed charge-polarization branch in the new completion direction \(H_{AB}\), the first unresolved term linear in \(H_{AB}\) is
\begin{equation}
\delta_H S_{\mathrm{unres}}\big|_{H=0}
=
\int_{\Sub} d^4X\,\sqrt{G}\,\bigl(\ResidualLift\bigr)^{AB}\,\delta H_{AB}.
\label{eq:linearcoupling}
\end{equation}
Equivalently, the minimal completion-source coupling is
\begin{equation}
\int_{\Sub} d^4X\,\sqrt{G}\,\bigl(\ResidualLift\bigr)^{AB}H_{AB}.
\label{eq:sourceaction}
\end{equation}
\end{proposition}

\begin{proof}
All lower-order bridge-response and charge-polarization contributions have already been absorbed into \(\BaseLift\) and \(\ResponseLift\). What remains unresolved at the next active burden is exactly the quartic self-response tensor \(\ResidualLift\) of Definition \ref{def:source}. Therefore the Taylor expansion of the still-unresolved action in a fresh completion direction \(H_{AB}\) starts with the linear contraction \eqref{eq:linearcoupling}. Any lower-order term vanishes by the earlier eliminations, and any higher completion self-interaction is beyond the present quartic-completion order bookkeeping.
\end{proof}

\begin{definition}[Explicit completion-sector action]
\label{def:action}
The completion-sector action is the auxiliary block
\begin{equation}
\CompAction[H,\GaugeAux,\ModeAux]
\equiv
\int_{\Sub} d^4X\,\sqrt{G}\,
\left[
\frac12 H^{AB}\SubEinLin(H)_{AB}
-\GaugeAux^A\SubGaugeVec(H)_A
+\frac12 \GaugeAux^A\GaugeAux_A
+(\ResidualLift)^{AB}H_{AB}
+\ModeAux^{AB}(\FreeProj H)_{AB}
\right].
\label{eq:Scmp}
\end{equation}
At the present note's scope, the completion sector is varied at fixed \((G,U,S,\ChargeFields,\psi)\). The branch-refined action is
\begin{equation}
\CompletedAction
\equiv
S_{\mathrm{cand}}\big|_{Q=(\ChargeFields,0,0,0)}
+\CompAction[H,\GaugeAux,\ModeAux].
\label{eq:Scandsharp}
\end{equation}
\end{definition}

\begin{remark}
Equation \eqref{eq:Scandsharp} is not yet a claim of full microscopic completion. It is an explicit action-level realization of the missing completion sector inside the already-fixed substrate action class. The universal matter route remains the same one-metric route of the underlying candidate branch and is reinstated after elimination through the completed metric defined below.
\end{remark}

\section{Elimination of the Constrained Auxiliary Sector}

\begin{proposition}[Eliminating the reduced-gauge auxiliary]
\label{prop:Y}
Variation of \(\CompAction\) with respect to \(\GaugeAux_A\) yields
\begin{equation}
\GaugeAux_A=\SubGaugeVec(H)_A.
\label{eq:Ysolve}
\end{equation}
After substituting \eqref{eq:Ysolve}, the completion action becomes
\begin{equation}
\CompAction[H,\GaugeAux(H),\ModeAux]
=
\int_{\Sub} d^4X\,\sqrt{G}\,
\left[
\frac12 H^{AB}\SubReducedEin(H)_{AB}
+(\ResidualLift)^{AB}H_{AB}
+\ModeAux^{AB}(\FreeProj H)_{AB}
\right].
\label{eq:reducedaction}
\end{equation}
\end{proposition}

\begin{proof}
Variation of \eqref{eq:Scmp} with respect to \(\GaugeAux_A\) gives
\[
\delta_{\GaugeAux}\CompAction
=
\int_{\Sub} d^4X\,\sqrt{G}\,
\delta \GaugeAux^A\bigl(-\SubGaugeVec(H)_A+\GaugeAux_A\bigr),
\]
so stationarity implies \eqref{eq:Ysolve}. Substituting \eqref{eq:Ysolve} back into \eqref{eq:Scmp} and integrating by parts in the standard reduced-harmonic way converts the \(-\GaugeAux\SubGaugeVec+\frac12 \GaugeAux^2\) block into \(-\frac12 H^{AB}\SubGaugeBook(H)_{AB}\), which yields \eqref{eq:reducedaction}.
\end{proof}

\begin{corollary}[Derived reduced completion operator]
\label{cor:derivedL}
The reduced completion operator entering the quartic-completion solve is
\[
\SubReducedEin=\SubEinLin-\SubGaugeBook,
\]
and it is therefore derived from the explicit auxiliary sector rather than postulated independently.
\end{corollary}

\begin{proposition}[Zero-extra-homogeneous-mode condition]
\label{prop:Z}
Variation of \(\CompAction\) with respect to \(\ModeAux_{AB}\) yields
\begin{equation}
\FreeProj H=0.
\label{eq:freemodezero}
\end{equation}
\end{proposition}

\begin{proof}
The multiplier \(\ModeAux_{AB}\) enters \eqref{eq:Scmp} linearly. Varying it gives
\[
\delta_{\ModeAux}\CompAction
=
\int_{\Sub} d^4X\,\sqrt{G}\,
\delta \ModeAux^{AB}(\FreeProj H)_{AB},
\]
hence \eqref{eq:freemodezero} at stationarity.
\end{proof}

\begin{theorem}[On-shell completion equation and recovery of the anchoring functional]
\label{thm:onshell}
Let \(\CompLift_{AB}\) denote the stationary completion tensor. Then
\begin{equation}
\SubReducedEin(\CompLift)_{AB}
=
-\ResidualLift_{AB},
\label{eq:subsolve}
\end{equation}
and
\begin{equation}
\FreeProj\CompLift=0.
\label{eq:nofree}
\end{equation}
Moreover, after elimination of \(\GaugeAux_A\) and \(\ModeAux_{AB}\), the completion action reduces exactly to the reconstruction-level functional of the anchoring note:
\begin{equation}
\CompFunctional[H]
\equiv
\frac12
\int_{\Sub} d^4X\,\sqrt{G}\,H^{AB}\SubReducedEin(H)_{AB}
+
\int_{\Sub} d^4X\,\sqrt{G}\,\bigl(\ResidualLift\bigr)^{AB}H_{AB},
\label{eq:functional}
\end{equation}
varied on the constrained class \(\FreeProj H=0\).
\end{theorem}

\begin{proof}
After eliminating \(\GaugeAux_A\) by Proposition \ref{prop:Y}, variation of \eqref{eq:reducedaction} with respect to \(H_{AB}\) gives
\begin{equation}
\SubReducedEin(H)_{AB}
+(\FreeProj^\ast \ModeAux)_{AB}
=
-\ResidualLift_{AB},
\label{eq:preproj}
\end{equation}
where \(\FreeProj^\ast=\FreeProj\) because \(\FreeProj\) is a projector on the admissible completion class. Applying \(\FreeProj\) to \eqref{eq:preproj} and using Definition \ref{def:freeproj} yields
\[
\FreeProj\SubReducedEin(H)
+\FreeProj\ModeAux
=
-\FreeProj\ResidualLift
=
0.
\]
Since \(\FreeProj\circ\SubReducedEin=0\), this gives \(\FreeProj\ModeAux=0\), hence \(\ModeAux=0\) on the multiplier image. Therefore \eqref{eq:preproj} reduces to \eqref{eq:subsolve}. Proposition \ref{prop:Z} gives \eqref{eq:nofree}. Finally, substituting \(\ModeAux=0\) into \eqref{eq:reducedaction} leaves exactly \eqref{eq:functional}, which is the anchoring functional.
\end{proof}

\begin{remark}
The action-level gain is now explicit. The source coupling \(\bigl(\ResidualLift\bigr)^{AB}H_{AB}\), the reduced operator \(\SubReducedEin\), and the zero-extra-homogeneous-mode condition \(\FreeProj H=0\) all arise from one constrained auxiliary sector. The earlier reconstruction-level functional is recovered as the eliminated form of that sector.
\end{remark}

\section{Observer Pullback and Recovery of the Same Completion Solve}

\begin{definition}[Completed bridge and observer metrics]
\label{def:completed}
Define the completed bridge metric and observer completion shift by
\begin{equation}
\BridgeMetric_{AB}^{\sharp}
\equiv
\BridgeMetric_{AB}
+\BaseLift_{AB}
+\ResponseLift_{AB}
+\CompLift_{AB},
\label{eq:bridgecompleted}
\end{equation}
\begin{equation}
\CompShift_{\mu\nu}
\equiv
\ObserverMap^*\CompLift_{AB},
\label{eq:defsigma}
\end{equation}
and
\begin{equation}
\CompletedMetric_{\mu\nu}
\equiv
\ObserverMetric_{\mu\nu}
+\BaseShift_{\mu\nu}
+\ResponseShift_{\mu\nu}
+\CompShift_{\mu\nu}.
\label{eq:gsharp}
\end{equation}
\end{definition}

\begin{theorem}[The same observer completion solve is recovered]
\label{thm:pullback}
The observer pullback of the on-shell substrate completion tensor satisfies exactly the same quartic-completion solve used in the redesign and anchoring notes:
\begin{equation}
\ReducedEin(\CompShift)_{\mu\nu}
=
-\ResidualCore{}_{\mu\nu}.
\label{eq:observereq}
\end{equation}
Equivalently,
\begin{equation}
\EinLin(\CompShift)_{\mu\nu}
=
-\ResidualCore{}_{\mu\nu}
+\GaugeBook(\CompShift)_{\mu\nu}.
\label{eq:linsolve}
\end{equation}
\end{theorem}

\begin{proof}
Apply \(\ObserverMap^*\) to the substrate equation \eqref{eq:subsolve}. Using \eqref{eq:intertwineL} and \eqref{eq:sourcepull},
\[
\ObserverMap^*\bigl(\SubReducedEin(\CompLift)\bigr)
=
\ReducedEin(\ObserverMap^*\CompLift)
=
\ReducedEin(\CompShift)
=
-\ObserverMap^*\ResidualLift
=
-\ResidualCore.
\]
This is \eqref{eq:observereq}. Equation \eqref{eq:linsolve} is the same identity rewritten using \(\ReducedEin=\EinLin-\GaugeBook\).
\end{proof}

\begin{corollary}[Equivalence with the earlier completion notes]
\label{cor:equiv}
Every observer-equation statement of the quartic-completion redesign note and the substrate-anchoring note applies verbatim to \(\CompletedMetric\), because the defining solve \eqref{eq:observereq} is unchanged.
\end{corollary}

\begin{proof}
Those notes use only the fixed candidate equation, the source decomposition \eqref{eq:vres}, and the completion solve \(\ReducedEin(\CompShift)=-\ResidualCore\). Theorem \ref{thm:pullback} re-establishes that same solve from the explicit action-level completion sector.
\end{proof}

\section{Exact Observer Equation on the Action-Derived Completed Branch}

\begin{definition}[Completion-generated matter and Einstein terms]
\label{def:extra}
Define
\begin{equation}
\DeltaTComp
\equiv
T_{\mu\nu}^{\mathrm{matter}}[\CompletedMetric,\psi]
-T_{\mu\nu}^{\mathrm{matter}}[\CandidateMetric,\psi],
\label{eq:deltatcomp}
\end{equation}
\begin{align}
\CompMix
&\equiv
\EinQuad[\BaseShift+\ResponseShift+\CompShift]
-\EinQuad[\BaseShift+\ResponseShift]
-\EinQuad[\CompShift],
\label{eq:compmix}
\\
\CompQuad
&\equiv
\EinQuad[\CompShift],
\label{eq:compquad}
\end{align}
and
\begin{equation}
\CompDec
\equiv
\ResidualDec+\CompMix+\CompQuad-8\pi G_N\,\DeltaTComp.
\label{eq:compdec}
\end{equation}
\end{definition}

\begin{theorem}[Exact observer equation on the action-derived completed branch]
\label{thm:obs}
The completed observer metric satisfies
\begin{equation}
G_{\mu\nu}[\CompletedMetric]
=
8\pi G_N\,T_{\mu\nu}^{\mathrm{matter}}[\CompletedMetric,\psi]
+\GaugeBook(\CompShift)_{\mu\nu}
+\CompDec.
\label{eq:newobs}
\end{equation}
In particular, the old genuine composite quartic residual tensor \(\ResidualCore\) is absent from the exact observer equation.
\end{theorem}

\begin{proof}
As in the earlier completion notes, expand the Einstein tensor about \(\CandidateMetric\):
\begin{equation}
G_{\mu\nu}[\CompletedMetric]
=
G_{\mu\nu}[\CandidateMetric]
+\EinLin(\CompShift)_{\mu\nu}
+\CompMix+\CompQuad.
\label{eq:expand}
\end{equation}
Substitute the fixed candidate equation
\[
G_{\mu\nu}[\CandidateMetric]
=
8\pi G_N\,T_{\mu\nu}^{\mathrm{matter}}[\CandidateMetric,\psi]
+\ResidualCore+\ResidualDec
\]
together with \eqref{eq:linsolve}. The \(\ResidualCore\) terms cancel exactly, leaving
\[
G_{\mu\nu}[\CompletedMetric]
=
8\pi G_N\,T_{\mu\nu}^{\mathrm{matter}}[\CandidateMetric,\psi]
+\GaugeBook(\CompShift)_{\mu\nu}
+\ResidualDec+\CompMix+\CompQuad.
\]
Replacing the candidate matter tensor by
\[
T_{\mu\nu}^{\mathrm{matter}}[\CandidateMetric,\psi]
=
T_{\mu\nu}^{\mathrm{matter}}[\CompletedMetric,\psi]-\DeltaTComp
\]
gives \eqref{eq:newobs}.
\end{proof}

\begin{remark}
The observer-equation gain is therefore unchanged but now sits one layer deeper. The quartic residual cancellation is no longer only a redesign choice and no longer only a reconstruction-level substrate functional. It is the observer pullback of an eliminated completion action inside the substrate action class.
\end{remark}

\section{Preservation of the One-Metric Route and the Recorded Phase Package}

\begin{proposition}[One operative metric is preserved]
\label{prop:onemetric}
On the action-derived completed branch, matter and light still couple only through the single operative metric \(\CompletedMetric\).
\end{proposition}

\begin{proof}
The underlying matter route of Definition \ref{def:candidate} remains \(\MatterAct[\ObserverMap^*\BridgeMetric,\psi]\). After completion, the bridge tensor entering that same route is \(\BridgeMetric^\sharp\), whose observer pullback is \(\CompletedMetric\) by Definition \ref{def:completed}. No second matter metric and no direct unsuppressed coupling to \((U^A,S,\ChargeFields,\GaugeAux_A,\ModeAux_{AB})\) is introduced.
\end{proof}

\begin{proposition}[Phase D / E are preserved]
\label{prop:phaseDE}
The action-derived completion preserves the recorded Phase D infrared-spectrum verdict and the Phase E universal-coupling verdict.
\end{proposition}

\begin{proof}
Linearizing \eqref{eq:subsolve} and \eqref{eq:nofree} on the admissible homogeneous benchmark background gives
\[
\SubReducedEin(\delta\CompLift)=0,
\qquad
\FreeProj(\delta\CompLift)=0.
\]
The only admissible free homogeneous completion mode is the one extracted by \(\FreeProj\), so \(\delta\CompLift=0\), hence
\begin{equation}
\delta\CompShift_{\mu\nu}=0.
\label{eq:nolinear}
\end{equation}
Therefore the linearized observer equation is unchanged relative to the already-screened charge-polarization candidate. No new unsuppressed scalar, vector, ghost, tachyonic, preferred-frame, or second-cone pole appears, and Proposition \ref{prop:onemetric} preserves the same single-metric universal matter route. That is exactly the Phase D / E package.
\end{proof}

\begin{proposition}[Infrared size of the action-derived completion]
\label{prop:size}
On every controlled infrared family,
\begin{equation}
\CompShift=\mathcal O(\epsIR^4),
\qquad
\GaugeBook(\CompShift)=\mathcal O(\epsIR^4),
\label{eq:ircomp}
\end{equation}
and
\begin{equation}
\CompMix=\mathcal O(\epsIR^6),
\qquad
\CompQuad=\mathcal O(\epsIR^8),
\qquad
\DeltaTComp=\mathcal O(\epsIR^4),
\qquad
\CompDec=\mathcal O(\epsIR^4).
\label{eq:irsizes}
\end{equation}
\end{proposition}

\begin{proof}
By \eqref{eq:sourceIR}, the source satisfies \(\ResidualLift=\mathcal O(\epsIR^4)\). The completion solve \eqref{eq:subsolve} is the same reduced solve used in the redesign and anchoring notes and is imposed with the same forced Coulomb-plus-regular class and zero free homogeneous mode \eqref{eq:nofree}; hence \(\CompLift=\mathcal O(\epsIR^4)\), and therefore \(\CompShift=\mathcal O(\epsIR^4)\). Since \(\GaugeBook\) is linear in first derivatives of \(\CompShift\), \eqref{eq:ircomp} follows. The remaining orders then follow exactly as in the anchoring note from
\[
\BaseShift=\mathcal O(\epsIR^2),
\qquad
\ResponseShift=\mathcal O(\epsIR^4),
\qquad
\CompShift=\mathcal O(\epsIR^4).
\]
\end{proof}

\begin{theorem}[Phase F controlled GR limit is preserved]
\label{thm:phaseF}
On every controlled infrared family,
\begin{equation}
G_{\mu\nu}[\CompletedMetric_{\epsIR}]
=
8\pi G_N\,T_{\mu\nu}^{\mathrm{matter}}[\CompletedMetric_{\epsIR},\psi_{\epsIR}]
+\GaugeBook(\CompShift_{\epsIR})_{\mu\nu}
+\epsIR^4\mathfrak D_{\mu\nu}^{(\epsIR)},
\label{eq:phaseFnew}
\end{equation}
with \(\mathfrak D_{\mu\nu}^{(\epsIR)}\) uniformly bounded. Thus the same one-metric observer equation reduces to Einstein plus ordinary matter as \(\epsIR\to 0\).
\end{theorem}

\begin{proof}
The exact observer equation \eqref{eq:newobs} holds on the completed branch. Proposition \ref{prop:size} gives \(\CompDec=\mathcal O(\epsIR^4)\), so \(\CompDec=\epsIR^4\mathfrak D^{(\epsIR)}\) with \(\mathfrak D^{(\epsIR)}\) uniformly bounded. The explicit remainder \(\GaugeBook(\CompShift)\) is still gauge bookkeeping rather than a genuine observer-side deviation sector, and Proposition \ref{prop:onemetric} preserves the same one operative metric. Therefore the same controlled GR-limit conclusion survives unchanged.
\end{proof}

\section{Claim-Boundary Verdict}

\begin{theorem}[Primary action-level completion verdict]
\label{thm:verdict}
At the disciplined scope of the present note, the positive result is exactly the following.
\begin{enumerate}
    \item The reconstruction-level completion functional of the anchoring note is derived from an explicit completion-sector action inside the substrate action class \eqref{eq:Scand}.
    \item The quartic self-response source coupling \(\bigl(\ResidualLift\bigr)^{AB}H_{AB}\) is derived as the first unresolved branch coupling in the completion direction.
    \item The reduced completion operator \(\SubReducedEin=\SubEinLin-\SubGaugeBook\) is derived by eliminating the auxiliary covector \(\GaugeAux_A\).
    \item The zero-extra-homogeneous-mode condition \(\FreeProj H=0\) is derived by eliminating the multiplier \(\ModeAux_{AB}\).
    \item Elimination of the action-derived sector yields the same substrate solve \(\SubReducedEin(\CompLift)=-\ResidualLift\), hence the same observer solve \(\ReducedEin(\CompShift)=-\ResidualCore\).
    \item The exact observer-equation closure, the one operative metric, and the recorded Phase D / E and Phase F package are preserved.
    \item Stronger promotion language is still not justified here, because the note gives an explicit action-level completion origin inside the existing branch but does not yet claim a unique microscopic ultraviolet derivation or a broader theorem than the current claim boundary allows.
\end{enumerate}
\end{theorem}

\begin{proof}
Item (1) is Definition \ref{def:action} and Theorem \ref{thm:onshell}. Item (2) is Proposition \ref{prop:sourcecoupling}. Item (3) is Proposition \ref{prop:Y} and Corollary \ref{cor:derivedL}. Item (4) is Proposition \ref{prop:Z}. Item (5) is Theorem \ref{thm:onshell} together with Theorem \ref{thm:pullback}. Item (6) is Theorem \ref{thm:obs}, Proposition \ref{prop:onemetric}, Proposition \ref{prop:phaseDE}, and Theorem \ref{thm:phaseF}. Item (7) is the claim-boundary discipline stated in the introduction.
\end{proof}

\begin{corollary}[What remains next]
The active next burden is no longer to invent another observer-level packaging move. It is to deepen the microscopic origin of the completion kernel itself, or else to produce another comparably concrete observer-equation gain on the same one-metric GR line. The anchored positive-pair line remains side work unless it contributes exactly that.
\end{corollary}

\section{Conclusion}

The substrate-anchoring note already showed that the quartic-completion solve could be viewed as the elimination of a reconstruction-level substrate functional. The present manuscript takes the next disciplined step below that result. Starting from the substrate action class already fixed by the derivational-bridge manuscript, it refines the extra substrate block \(Q\) by adjoining an explicit completion tensor, a reduced-gauge auxiliary covector, and a free-homogeneous-mode multiplier. From that action-level sector it derives the quartic source coupling, the reduced completion operator, and the zero-extra-homogeneous-mode condition.

Eliminating that constrained auxiliary sector reproduces exactly the same completion equation already used on the active charge-polarization line:
\[
\SubReducedEin(\CompLift)=-\ResidualLift,
\qquad
\ReducedEin(\CompShift)=-\ResidualCore.
\]
Consequently, the completed observer equation is again free of the old genuine composite quartic residual tensor, the one operative metric is preserved, and the recorded Phase D / E and Phase F package remains intact.

The scope of the result remains honest. This note does not yet claim a unique microscopic ultraviolet completion, does not widen the current finite-amplitude theorem boundary, and does not license stronger promotion language. Its positive result is narrower and exactly the next one needed: the reconstruction-level quartic-completion functional now has an explicit action-level origin inside the same substrate action class and the same one-metric charge-polarization branch.

\end{document}
